% =============================================================================
% Section 2 — Methods
% =============================================================================
\section{Methods}
\label{sec:methods}

\subsection{Abdominal shell model}
\label{sec:shell-model}

We adopt the fluid-filled viscoelastic oblate spheroidal shell model
developed in our companion paper~\cite{Mace2026browntone}.  The abdomen
is represented as an oblate spheroid with semi-major axis
$a = \SI{0.18}{\metre}$, semi-minor axis $c = \SI{0.12}{\metre}$,
uniform wall thickness $h = \SI{0.010}{\metre}$, and equivalent sphere
radius $R_\mathrm{eq} = (a^2 c)^{1/3} = \SI{0.157}{\metre}$.  The wall
is characterised by Young's modulus $E = \SI{0.1}{\mega\pascal}$,
Poisson's ratio $\nu = 0.45$, density
$\rho_w = \SI{1100}{\kilogram\per\metre\cubed}$, and loss tangent
$\eta = 0.25$ (quality factor $Q = 4$, damping ratio $\zeta = 0.125$).
The enclosed fluid has density
$\rho_f = \SI{1020}{\kilogram\per\metre\cubed}$ and bulk modulus
$K_f = \SI{2.2}{\giga\pascal}$; the intra-abdominal pressure is
$P_\mathrm{iap} = \SI{1000}{\pascal}$.

The flexural mode natural frequencies follow from the Rayleigh quotient
for a pressurised, fluid-loaded spherical shell~\cite{Junger2012, Lamb1882}:
%
\begin{equation}
  \omega_n^2 = \frac{K_\mathrm{bend}(n) + K_\mathrm{memb}(n)
    + K_\mathrm{pre}(n)}{m_\mathrm{wall} + m_\mathrm{added}(n)}
  \label{eq:flexural-freq}
\end{equation}
%
where the bending, membrane, and pre-stress stiffnesses per unit area
are:
%
\begin{align}
  K_\mathrm{bend} &= n(n-1)(n+2)^2 \frac{D}{R^4}, \\
  K_\mathrm{memb} &= \frac{Eh}{R^2}\,
    \frac{n^2+n-2+2\nu}{n^2+n+1-\nu}, \\
  K_\mathrm{pre}  &= \frac{P_\mathrm{iap}}{R}\,(n-1)(n+2),
\end{align}
%
and the effective mass per unit area is
$m_\mathrm{eff} = \rho_w h + \rho_f R / n$.  With canonical parameters
the $n = 2$ flexural resonance lies at
$f_2 \approx \SI{3.95}{\hertz}$---well below the sub-bass band.

\subsection{Acoustic coupling at sub-bass frequencies}
\label{sec:coupling}

For a plane wave incident on a sphere of radius $R$, the effective
driving pressure for flexural mode $n$ in the long-wavelength limit
is~\cite{Junger2012}:
%
\begin{equation}
  p_\mathrm{eff}(n) = p_\mathrm{inc} \, (ka)^n,
  \label{eq:ka-coupling}
\end{equation}
%
where $k = 2\pi f / c_\mathrm{air}$ is the acoustic wavenumber and
$a \equiv R_\mathrm{eq}$.  The Helmholtz number $ka$ is the
dimensionless parameter governing the transition from Rayleigh
scattering ($ka \ll 1$) to geometric acoustics ($ka \gg 1$).

At the $n = 2$ resonance (\SI{3.95}{\hertz}), $ka = 0.011$ and
$(ka)^2 \approx 1.3 \times 10^{-4}$, yielding negligible airborne
coupling~\cite{Mace2026browntone}.  In the sub-bass band:
%
\begin{equation}
  ka(f) = \frac{2\pi f R_\mathrm{eq}}{c_\mathrm{air}},
  \label{eq:ka-def}
\end{equation}
%
giving $ka \approx 0.057$ at \SI{20}{\hertz} and $ka \approx 0.23$ at
\SI{80}{\hertz}.  The corresponding $(ka)^2$ values are
$3.3 \times 10^{-3}$ and $5.3 \times 10^{-2}$---up to $\sim$400 times
the value at resonance.

\subsection{Off-resonance forced response}
\label{sec:sdof}

Although acoustic coupling improves with frequency, the driving
frequency is now far from the $n = 2$ resonance.  The response is
governed by the single-degree-of-freedom (SDOF) transfer function:
%
\begin{equation}
  H(r, \zeta) = \frac{1}{\sqrt{(1 - r^2)^2 + (2\zeta r)^2}},
  \label{eq:sdof}
\end{equation}
%
where $r = f / f_n$ is the frequency ratio and $\zeta$ is the damping
ratio.  At resonance ($r = 1$), $H = Q = 1/(2\zeta) = 4$.  At
$f = \SI{40}{\hertz}$ ($r \approx 10$), $H \approx 1/r^2 \approx 0.01$.

The tissue displacement at frequency $f$ and sound pressure level
$L_p$~dB is:
%
\begin{equation}
  \xi(f) = \frac{p_\mathrm{eff}(f)}{K_\mathrm{total}} \, H(r, \zeta)
  = \frac{p_\mathrm{inc}(L_p)\,(ka)^n}{K_\mathrm{total}} \, H(r, \zeta).
  \label{eq:displacement}
\end{equation}

The net effect is a competition: $(ka)^n$ increases with frequency
(improving coupling), but $H(r)$ decreases (worsening off-resonance
attenuation).  The question is which dominates.

\subsection{Concert sound pressure level spectra}
\label{sec:concert-spl}

We use literature-derived third-octave spectra for three concert
environments:
%
\begin{itemize}
  \item \textbf{Rock concert} (broadband, \SIrange{95}{108}{\decibel}
    sub-bass)~\cite{Gunderson2016}.
  \item \textbf{Electronic dance music} (sub-woofer-driven, peak
    \SI{115}{\decibel} at \SIrange{31.5}{50}{\hertz})~\cite{Opperman2006}.
  \item \textbf{Pipe organ} (extreme low-frequency extension, 32-foot
    rank fundamental at \SI{16}{\hertz}, peak \SI{105}{\decibel} at
    \SIrange{31.5}{50}{\hertz})~\cite{Bickerdike1966}.
\end{itemize}

For each genre, the SPL at each frequency is interpolated from published
third-octave measurements and used as the input to
equation~\eqref{eq:displacement}.

\subsection{Perception thresholds}
\label{sec:perception}

We compare computed tissue displacements against whole-body vibration
perception thresholds from ISO~2631-1:1997~\cite{iso2631}.  The standard
specifies root-mean-square acceleration thresholds as a function of
frequency; we convert these to displacement amplitudes via:
%
\begin{equation}
  \xi_\mathrm{thresh}(f) = \frac{a_\mathrm{rms}(f)\,\sqrt{2}}{(2\pi f)^2}.
  \label{eq:perception-disp}
\end{equation}

This provides a conservative comparison: ISO~2631 thresholds are for
whole-body mechanical vibration, not for localised tissue displacement.
If the airborne-induced displacement is below even these thresholds, it
is certainly imperceptible.
